\subsection{Comparing candidates under shared uncertainty}
\label{app:paired-candidate-bounds}

Separate score intervals discard the dependence induced by a common actual
state. For example, costs $10+3\upsilon$ and $12+3\upsilon$, with
$\upsilon\in[-1,1]$, have overlapping intervals but differ by exactly $-2$.
This section applies dependence-preserving range analysis
\citep{stolfi2003affine} to the assignments and rewiring in a candidate query.
Absolute feasibility remains governed by Corollary~\ref{cor:query-certificate}.

\paragraph{Joint uncertainty after assignment.}
Fix the finite candidate family of Algorithm~\ref{alg:cphast_query} and a
joint physical uncertainty box, normalized as $\upsilon\in[-1,1]^m$.
Each coordinate represents one uncertain physical quantity, reused across
candidates only when it is the same quantity. Candidate-specific additions
receive distinct coordinates. The box is a joint set, not an assumption of
stochastic independence. For affine coordinate conversions and resets,
gathering retained states and inserting supplied new states gives
\[
 x_0^a(\upsilon)=d_a+D_a\upsilon,
 \qquad d_a\in\mathbb R^{n_a},\quad D_a\in\mathbb R^{n_a\times m}.
\]
The center and radii are included in $d_a,D_a$; $n_a$ may differ between
candidates. Assume exact known affine sampled dynamics
$x_{k+1}^a=F_{a,k}x_k^a+f_{a,k}$ over the fixed horizon, with prescribed inputs
included in $f_{a,k}$ and all trajectories in the region of validity.
Stacked affine readouts then have the form
$r_a(\upsilon)=\bar r_a+T_a\upsilon\in\mathbb R^{p_a}$.
For a fixed quadratic physical score, abbreviate
\begin{equation}
 J_a(\upsilon):=J_w(a;\upsilon)
 =r_a^\top W_a r_a+v_a^\top r_a+c_{0,a},
 \qquad W_a=W_a^\top\in\mathbb R^{p_a\times p_a},\quad
 v_a\in\mathbb R^{p_a}.
 \label{eq:paired-score-setup}
\end{equation}
All stage weights, normalizations, and terminal terms are fixed in this
score. Peak-current constraints are screened separately, not treated as
quadratic objectives.

\begin{proposition}[Assignment-aware candidate-difference enclosure]
\label{prop:paired-candidate-enclosure}
Under the preceding assumptions, each score is exactly
$J_a=c_a+\ell_a^\top\upsilon+\upsilon^\top Q_a\upsilon$, where
\begin{equation}
 \begin{aligned}
 Q_a&=T_a^\top W_a T_a,\qquad
 \ell_a=T_a^\top(2W_a\bar r_a+v_a),\\
 c_a&=\bar r_a^\top W_a\bar r_a+v_a^\top\bar r_a+c_{0,a}.
 \end{aligned}
 \label{eq:paired-score-coefficients}
\end{equation}
For a quadratic $p(\upsilon)=c+\ell^\top\upsilon+\upsilon^\top Q\upsilon$
with $Q=Q^\top$, define the monomial enclosure
\begin{equation}
 \begin{aligned}
 t(p)&=\sum_i|\ell_i|+2\sum_{i<j}|Q_{ij}|,\\
 \mathcal E(p)&=
 \left[c+\sum_i\min(0,Q_{ii})-t(p),\;
       c+\sum_i\max(0,Q_{ii})+t(p)\right].
 \end{aligned}
 \label{eq:paired-monomial-enclosure}
\end{equation}
Then the range of $J_a-J_b$ is contained in
$\mathcal E(J_a-J_b)$, computed by subtracting the coefficients in
Eq.~\eqref{eq:paired-score-coefficients}. With interval subtraction on the right,
\begin{equation}
 \mathcal E(J_a-J_b)\subseteq\mathcal E(J_a)-\mathcal E(J_b),
 \label{eq:paired-enclosure-inclusion}
\end{equation}
\end{proposition}

\begin{proof}
Affine propagation preserves the shared coordinates: if
$x_k^a=t_k+T_k\upsilon$, then
$t_{k+1}=F_{a,k}t_k+f_{a,k}$ and $T_{k+1}=F_{a,k}T_k$.
Stacking affine readouts and expanding Eq.~\eqref{eq:paired-score-setup}
gives Eq.~\eqref{eq:paired-score-coefficients}.
Each linear and off-diagonal monomial ranges in $[-1,1]$, whereas
$\upsilon_i^2\in[0,1]$; summing their signed ranges gives
Eq.~\eqref{eq:paired-monomial-enclosure}.
For any common monomial range $I$ and coefficients $\alpha,\beta$,
$(\alpha-\beta)I\subseteq\alpha I-\beta I$, since the right side permits
different values of that monomial. Summing these inclusions over the same
constant, linear, diagonal, and off-diagonal basis proves
Eq.~\eqref{eq:paired-enclosure-inclusion}.
\end{proof}

A strictly negative upper endpoint certifies $J_a<J_b$ throughout the joint
box; a strictly positive lower endpoint certifies the reverse.
The inclusion compares the \emph{same monomial representation}; a factored
expression or tighter marginal optimizer can do better. For any valid
individual score intervals $[L_a,U_a]$ and $[L_b,U_b]$, use
\begin{equation}
 \mathcal I_{ab}=
 \mathcal E(J_a-J_b)\cap[L_a-U_b,U_a-L_b].
 \label{eq:paired-hybrid}
\end{equation}
Both intervals contain the true difference, so their intersection is valid
and never wider than either. Containing zero leaves the ordering unresolved;
an empty intersection indicates inconsistent enclosures. For any retained
choice $a$, valid paired upper endpoints $U_{ab}^{\mathrm{pair}}$, with
$U_{aa}^{\mathrm{pair}}=0$, also give
\begin{equation}
 J_a(\upsilon)-\min_{b\in\widehat{\mathcal A}_{\mathrm{safe}}}J_b(\upsilon)
 \le\max_{b\in\widehat{\mathcal A}_{\mathrm{safe}}}U_{ab}^{\mathrm{pair}},
 \label{eq:paired-retained-regret}
\end{equation}
because the left side equals $\max_b(J_a-J_b)$ over that same set.

\paragraph{Rewiring and residual errors.}
The quadratic term cancels only when $Q_a=Q_b$. Even with shared initial
states and local laws, rewiring changes the response matrices $T_a,T_b$;
the term $\upsilon^\top(Q_a-Q_b)\upsilon$ must therefore be retained.
Shared fitted weights alone do not imply shared physical model error.
For surrogate quadratic scores with certified residuals
$|J_a-p_a|\le\rho_a$, a valid difference enclosure is
$\mathcal E(p_a-p_b)+[-\rho_a-\rho_b,\rho_a+\rho_b]$.
Nonlinear assignment, uncertain-parameter flow, and learned-law errors
require such validated remainders; they are absent from the exact known-law
case below. Numerical coefficients must also be outward enclosed through
flow construction, coefficient assembly, and interval evaluation.

\paragraph{Known-law rewiring check.}
We use the three-DGU physical laws, initial state, nine source-voltage choices,
quadratic regulation/current-sharing score, and 100-step horizon of
Appendix~\ref{app:certificate-validation}. All resistances are exactly
$1\,\Omega$. Two prospective graphs use incidence matrices
\[
 B_{\mathrm{inc}}^{(1)}=
 \begin{bmatrix}1&0\\-1&1\\0&-1\end{bmatrix},\qquad
 B_{\mathrm{inc}}^{(2)}=
 \begin{bmatrix}1&1\\-1&0\\0&-1\end{bmatrix}.
\]
The second graph moves one line's positive endpoint. Both graphs retain the
same components and eight uncertain initial voltage/current coordinates,
with identity assignment and reset. Each graph remains fixed during its
counterfactual rollout; no switching transient is modeled. The 18 candidates
yield 81 cross-graph comparisons before screening. We use 50-decimal
outward-rounded interval arithmetic and the degree-40 exponential enclosure
of Appendix~\ref{app:certificate-validation}, preserving the signed
initial-state dependence through the quadratic coefficient construction.

\begin{table}[H]
\centering
\small
\caption{\textbf{Shared uncertainty resolves additional rewiring comparisons.}
Number of strictly certified cross-graph orderings after the same synthetic
2 A sampled-current screen. Radius $r$ is in V for voltage coordinates and
A for current coordinates. Separate raw intervals evaluate the reference
score from state boxes; separate polynomial intervals subtract
$\mathcal E(J_a),\mathcal E(J_b)$; paired bounds use
Eq.~\eqref{eq:paired-hybrid}. Bold green marks the largest count among
nonempty comparisons, including ties. Zero eligible pairs means no admitted
comparison, not successful certification.}
\label{tab:paired-rewiring-bounds}
\setlength{\tabcolsep}{7pt}
\begin{tabular}{crrrrr}
\toprule
Radius $r$ & Retained & Eligible pairs & Separate raw & Separate polynomial & Paired \\
\midrule
0     & 18 & 81 & \goodcell{\textbf{81}} & \goodcell{\textbf{81}} & \goodcell{\textbf{81}} \\
0.001 & 18 & 81 & \goodcell{\textbf{81}} & \goodcell{\textbf{81}} & \goodcell{\textbf{81}} \\
0.01  & 18 & 81 & 72 & 75 & \goodcell{\textbf{79}} \\
0.05  & 18 & 81 & 56 & 62 & \goodcell{\textbf{74}} \\
0.2   & 18 & 81 & 6  & 32 & \goodcell{\textbf{60}} \\
1     & 0  & 0  & 0  & 0  & 0 \\
\bottomrule
\end{tabular}
\end{table}

At $r=0.05$, pairing resolves 12 more cross-graph comparisons than separate
polynomial bounds, beyond the improvement from changing the marginal
representation (Table~\ref{tab:paired-rewiring-bounds}). Direct and intersected
counts coincide in this example.

Absolute feasibility is unchanged: the
stricter 1 A screen retains only six chain candidates through $r=0.05$, so
admits no cross-graph comparison; at $r\ge0.2$ it retains none. The 0.25 and
0.5 A screens retain none at any tested radius. At $r=1$, none of the four
screens admits a candidate.

The interval construction establishes the box-wide comparisons, conditional
on the known laws and correct outward arithmetic. Twelve exact-rational
regression tests and 756 finite reference paths check the implementation,
but are not a substitute for those assumptions. This is an
architecture-independent refinement of candidate comparison, not a
trained-model superiority or learned-observer coverage result. Its baseline
is the raw reference-score interval, not the predictor-centered envelope
used in Appendix~\ref{app:certificate-validation}.
\FloatBarrier
